% Exercise 2.8
To calculate the IRR on an incremental basis, we subtract each cash flow of
the less expensive of the two options considered from the more expensive
option, yielding positive cash flows for all but the initial payment. Thus we
have the following table of cash flows:
\begin{center}
    \begin{tabular}{l|r|r|r}
        & Options A vs B & Options B vs C & Options A vs C \\
        \hline
        $x_0$ & -24,000 & -5,000 & -29,000 \\
        $x_1 \cdots x_4$ & 6,000 & 400 & 6,400 \\
        $x_5$ & 10,000 & 2,000 & 12,000
    \end{tabular}
\end{center}

Thus we have the following polynomial for considering changing from Option A
to Option B:
\[
    -24,000 + 6,000c + 6,000c^2 + 6,000c^3 + 6,000c^4 + 10,000c^5,
\]
or, equivalently
\[
    p(r) = -24,000 + \frac{6,000}{1 + r} + \frac{6,000}{{(1 + r)}^2}
    + \frac{6,000}{{(1 + r)}^3} + \frac{6,000}{{(1 + r)}^4}
    + \frac{10,000}{{(1 + r)}^5}.
\]

Using Newton's method from Exercise 4, we have the following calculation:
\begin{lstlisting}
>>> from investment_science.cash_flows import irr_iterates, present_value
>>> from investment_science.ch02.ex2_8 import A_VS_B
>>> for r in irr_iterates(A_VS_B, 0.1, 5):
...     print(f"r = {r}, f(r) = {present_value(A_VS_B, r)}")
...
r = 0.1, f(r) = 1228.405908687304
r = 0.11770177473592071, f(r) = 45.505742594757976
r = 0.11840909697075934, f(r) = 0.06799356632063791
r = 0.11841015700002991, f(r) = 1.5231762517942116e-07
r = 0.11841015700240458, f(r) = 8.185452315956354e-12
r = 0.1184101570024047, f(r) = -6.821210263296962e-12
\end{lstlisting}
Thus we have $r_0 \approx 0.118$, i.e., an IRR of approximately 11.8\%,
indicating that a change from A to B is justified on the basis of the IRR.
